\begin{table*}[!tp]
\caption{Candidate transition families for resolving
hyperfine groups, with L26 predictions.}
%$J$ is the upper rotational quantum number, $N$ the number of resolved
%positive response features, and $\Delta\nu$ their adjacent peak
%separations. Ranges describe variation between transitions.}
\label{tab:experiment-families}
% BL review remains open: inspect the coupling-sensitivity interpretation; values unchanged.
\centering
\begingroup
\small
\setlength{\tabcolsep}{3pt}
\renewcommand{\arraystretch}{1.16}
\begin{ruledtabular}
\begin{tabular}{lrc lr}
Family & $\nu$ (GHz)\textsuperscript{a} & $J$\textsuperscript{d} & Coupling sensitivity\textsuperscript{b} & $N$; $\Delta\nu$ (kHz)\textsuperscript{c}\\
\hline
\multicolumn{5}{l}{$v_{\rm t}=0$}\\
$J_{-2}\leftarrow J_2\,E$\hspace{0.6em} & 118.523--239.659 & 24--33 & \shortstack[l]{Linked methyl SR;\\OH $M^{\rm s}_{ab,0},d^{\prime}_{b,0}$} & 4; 10.5--24.4 \\[3pt]
$J_1\leftarrow J_{-1}\,E$\hspace{0.6em} & 435.400--471.899 & 30--33 & \shortstack[l]{Linked methyl SR;\\OH $M^{\rm s}_{ab,0},d^{\prime}_{b,0}$} & 4; 13.9--33.7 \\[3pt]
$J_1\leftarrow J_0\,E$\hspace{0.6em} & 259.581--312.995 & 24--28 & \shortstack[l]{OH $M^{\rm s}_{ab,0},d^{\prime}_{b,0}$;\\methyl SR} & 4; 8.9--20.5 \\[3pt]
$J_2\leftarrow J_{-1}\,E$\hspace{0.6em} & 368.731--491.769 & 17--34 & \shortstack[l]{Methyl SR, $n=2$; finer linked\\SR and OH $M^{\rm s}_{ab,0},d^{\prime}_{b,0}$} & 2--4; 4.2--38.3 \\[3pt]
$J_{-2}\leftarrow J_1\,E$\hspace{0.6em} & 147.397--248.971 & 24--33 & \shortstack[l]{Methyl SR, $n=2$; finer linked\\SR and OH $M^{\rm s}_{ab,0},d^{\prime}_{b,0}$} & 4; 8.3--19.8 \\[3pt]
$(J+1)_0\leftarrow J_0\,A^+$\hspace{0.6em} & 48.372--96.741 & 0--1 & \shortstack[l]{Methyl--methyl and\\methyl--OH DD} & 2--4; 11.4--14.9 \\[3pt]
$J_2\leftarrow J_1\,E$\hspace{0.6em} & 24.929--25.541 & 2--9 & \shortstack[l]{Methyl--OH DD; transverse\\ST and OH $M^{\rm s}_{ab,0}$} & 1--3; 6.5--10.6 \\[3pt]
$J_1A^-\leftarrow J_1A^+$\hspace{0.6em} & 2.503--17.513 & 2, 3, 6 & \shortstack[l]{Main separation: SR;\\finer structure: DD} & 1--2; 7.3--7.3 \\[3pt]
\hline
\multicolumn{5}{l}{$v_{\rm t}=1$}\\
$(J+1)_7\leftarrow J_6\,E$\hspace{0.6em} & 153.129--537.726 & 7--15 & \shortstack[l]{Axial $d^{\prime}_{a,0}$ and transverse\\$d^{\prime}_{b,0}$ of methyl and OH} & 4; 22.7--111.2 \\[3pt]
$(J+1)_2\leftarrow J_3\,E$\hspace{0.6em} & 21.147--413.353 & 7--15 & \shortstack[l]{Axial and transverse ST;\\little OH $M^{\rm s}_{ab,0}$ response} & 4; 11.3--51.3 \\[3pt]
$(J+1)_7\leftarrow J_8\,E$\hspace{0.6em} & 140.843--566.568 & 8--17 & \shortstack[l]{Axial ST; OH $d^{\prime}_{b,0},M^{\rm s}_{ab,0}$;\\methyl transverse ST at low $J$} & 2--5; 5.2--80.9 \\[3pt]
$J_6\leftarrow J_7\,E$\hspace{0.6em} & 232.423--233.156 & 7--15 & \shortstack[l]{Axial $d^{\prime}_{a,0}$ and transverse\\$d^{\prime}_{b,0}$ of methyl and OH} & 4; 24.6--107.4 \\[3pt]
$J_8\leftarrow J_7\,E$\hspace{0.6em} & 292.864--318.331 & 8--24 & \shortstack[l]{Axial ST; OH $d^{\prime}_{b,0},M^{\rm s}_{ab,0}$;\\methyl transverse ST in close pairs} & 2--4; 14.6--98.6 \\[3pt]
$J_3\leftarrow J_2\,E$\hspace{0.6em} & 355.965--365.047 & 3--16 & \shortstack[l]{Axial ST; transverse ST\\and methyl--OH DD} & 4; 11.3--118.7 \\[3pt]
\end{tabular}
\end{ruledtabular}
\endgroup
\par\vspace{3pt}
\begin{minipage}{\linewidth}
\footnotesize\raggedright
\textsuperscript{a)} Frequency range of transition family. \par
\textsuperscript{b)} SR, ST, and DD denote spin--rotation ($M^{\rm s}$), effective
    spin--torsion ($\bm d'$), and dipole--dipole interaction (tensors).
The linked methyl direction predominantly contains
$M^{{\rm s},{\rm CH_3}}_{ac,1}$, $M^{{\rm s},{\rm CH_3}}_{ab,1}$,
and $d^{\prime\rm CH_3}_{c,1}$; methyl SR at $n=2$ denotes the
$bb$, $cc$, and $bc$ terms.\par
\textsuperscript{c)} $N$: resolved positive-response features; $\Delta\nu$: adjacent peak separations. Ranges span transitions.\par
Second-harmonic responses use Lorentzian components of HWHM
5~kHz, sinusoidal frequency-modulation amplitude 3~kHz, and statistical
strengths, without crossovers or saturation-induced redistribution.
The response phase is chosen so an isolated component has a positive
central peak. Retained peaks have height at least $10\%$ and prominence at
least $5\%$ of the profile maximum. Peak positions are refined continuously;
counts and gaps agree on 0.02- and 0.01-kHz sampling grids. These are numerical
feature-selection criteria, not calibrated experimental resolution limits.
$N=1$ denotes a merged profile and is excluded from the gap ranges;
this occurs for $J=6$--9 in the 25-GHz family and $J=2,3$ in the
$A$ parity family.\par
\textsuperscript{d)} $J$: lower-state rotational quantum number.
\end{minipage}
\end{table*}
